% En esta seccion se realiza una investigación formal del tema y lo presenta bien fundamentado, haciendo referencia a las fuentes de información usando las normas APA.

\section{Introducción}
En esta tercera práctica del Laboratorio de Computación Gráfica el trabajo se centró en el modelado geométrico. En la Práctica 02 habíamos construido una casa, pero el interés estaba puesto en la cámara y en las matrices que determinan cómo se le observa; aquí ocurre lo contrario: la cámara se deja casi fija y lo que se estudia es la figura misma, es decir, cómo se decide qué vértices la componen y de qué manera se unen entre sí. Un modelo geométrico es la representación de las características geométricas de una entidad, concreta o abstracta, y se construye a partir de primitivas, esto es, de elementos más sencillos que el objeto que se quiere representar \cite{teodoro_manual03}.

Conviene distinguir dos nociones que en el lenguaje cotidiano se confunden. Una \emph{primitiva geométrica} es la figura que el usuario tiene en mente: un plano, un cubo, un círculo, una esfera. Una \emph{primitiva gráfica}, en cambio, es lo que la interfaz sabe dibujar realmente, que en el perfil \textit{core} de OpenGL se reduce a puntos, líneas y triángulos \cite{shreiner_opengl}. Entre una y otra hay siempre un paso de descomposición: el cubo no se envía a la tarjeta como «un cubo», sino como doce triángulos, y el círculo, que en el papel es una curva continua, se aproxima con un polígono de $n$ lados cuyo error disminuye conforme $n$ crece. Esa aproximación es la razón por la que el modelado geométrico se estudia como un problema de parametrización y no sólo como una lista de coordenadas \cite{hearn_baker}.

La estructura con la que se describen estas figuras es la malla poligonal, formada por un arreglo de vértices y un arreglo de índices que indica qué tercias de vértices forman cada triángulo \cite{marschner_shirley}. Separar ambos arreglos permite que un mismo vértice participe en varias caras sin duplicarse, lo cual reduce la memoria empleada y, sobre todo, la cantidad de datos que se transfieren a la unidad de procesamiento gráfico. En el proyecto base esa estructura está encapsulada en la clase \texttt{Mesh}, que crea el objeto de arreglo de vértices y los búferes correspondientes y dibuja el conjunto con \texttt{glDrawElements} \cite{learnopengl_hello}. Cada vértice lleva además un color expresado en el modo RGBA, en el que las componentes cromáticas y la de opacidad se manejan como valores flotantes acotados \cite{microsoft_rgba}; como el color se declara por vértice, la etapa de rasterización interpola los valores intermedios y produce los degradados que se aprecian en las capturas de este reporte.

Las tres actividades que se documentan aquí recorren esa idea en orden creciente de dificultad. El plano es el caso mínimo, con cuatro vértices y dos triángulos, y sirve para fijar la convención de que la figura se construye centrada en el origen. El cubo introduce la tercera dimensión y obliga a razonar por caras, comprobando que ninguna quede sin cubrir. El círculo introduce por fin la parametrización: sus vértices no se escriben a mano, sino que se calculan con las funciones trigonométricas evaluadas en incrementos regulares del ángulo, de modo que el número de lados se vuelve un parámetro del constructor. Para cada figura se asignaron dos vistas a las teclas \texttt{F1} y \texttt{F2}, tal como pide el manual, y se agregó una tecla que alterna el modo de rasterización entre relleno y alambre mediante \texttt{glPolygonMode}, con la que puede verse directamente la descomposición en triángulos \cite{khronos_polygonmode}.

El desarrollo se llevó a cabo en C++ sobre el proyecto base proporcionado por el profesor, que emplea OpenGL 3.3 con GLFW para la gestión de la ventana, GLAD para la carga de las funciones de la API y GLM para las operaciones con matrices y vectores. La compilación se realizó en Visual Studio siguiendo la configuración descrita en el manual de la práctica. Durante el desarrollo se utilizó además un modelo de lenguaje como apoyo para analizar el comportamiento de la clase \texttt{Circle} del proyecto base y para la revisión de redacción del documento \cite{anthropic2026}.
